\documentclass[11pt]{article}
\usepackage[margin=25mm]{geometry}
\usepackage{fontspec}
\setmainfont{Latin Modern Roman}
\usepackage{amsmath,amssymb,mathtools}
\usepackage{unicode-math}
\setmathfont{Latin Modern Math}
\usepackage{xurl,hyperref}
\hypersetup{colorlinks=true,urlcolor=blue}
\setcounter{secnumdepth}{0}
\setlength{\parindent}{0pt}
\setlength{\parskip}{5pt}
\setlength{\emergencystretch}{3em}
\newcommand{\sref}[2]{\hyperlink{source:#1:#2}{[S#2]}}
\newcommand{\cataloguescope}[1]{\paragraph{Recorded target.}#1}
\begin{document}
\section{Research Problems in Function Theory — Problem 6.108}\label{statement}

\subsection{Definitions and mathematical statement}

\paragraph{Shared notation.}$\mathbb N=\{1,2,\ldots\}$, and $\mathbb Z,\mathbb Q,\mathbb R,\mathbb C$ denote the integers, rationals, reals and complex numbers. Any other conventions are specified locally.


Write $\mathbb D=\{z\in\mathbb C:|z|<1\}$ and $\mathbb T=\partial\mathbb D$. Let $S$ be the class of holomorphic injective functions $f:\mathbb D\to\mathbb C$ with $f(0)=0$ and $f'(0)=1$. For a holomorphic univalent $f$ in $\mathbb D$, $0<r<1$ and $\lambda>0$, let
\[I_\lambda(r,f')=\left(\frac1{2\pi}\int_0^{2\pi}|f'(re^{i\theta})|^\lambda\,d\theta\right)^{1/\lambda}.\]
At $\lambda=0$ use the limiting geometric mean
\[I_0(r,f')=\exp\left(\frac1{2\pi}\int_0^{2\pi}\log|f'(re^{i\theta})|\,d\theta\right).\]
The derivative never vanishes. For normalized $f\in S$, put $k(z)=z/(1-z)^2$ and
\[C_\lambda=\sup_{f\in S,\ 0<r<1}\frac{I_\lambda(r,f')}{I_\lambda(r,k')}.\]
\paragraph{Statement recorded in the supplied source.}
Determine the maximal order of growth of $I_\lambda(r,f')$ as $r\uparrow1$ for $0\leq\lambda\leq2/5$, as $f$ ranges over univalent functions. Bounds for an individual unnormalized $f$ may depend on $f$. For $\lambda>2/5$, determine precisely when $C_\lambda=1$; in particular, is $C_\lambda=1$ for every $\lambda\geq2$? \sref{statement}{2}
\subsection{Short English statement}
Find the small-exponent growth of derivative integral means, and determine when the Koebe function gives the exact normalized bound at larger exponents.
\subsection{Sources}
\begin{description}
\item[\hypertarget{source:statement:1}{S1}] ULAM.AI and the UnsolvedMath Contributors, UnsolvedMath: A Curated Collection of Open Mathematics Problems, catalogue entry. CC BY 4.0. \url{https://huggingface.co/datasets/ulamai/UnsolvedMath}.
\item[\hypertarget{source:statement:2}{S2}] W. K. Hayman and E. F. Lingham, Research Problems in Function Theory (2018), Problem 6.108, complete problem and update. \url{https://arxiv.org/abs/1809.07200}.
\end{description}
\label{end:statement}
\end{document}
